% GENERATED FILE. Edit canonical lessons, then run npm run generate:books.
% canonical-source-hash: fnv1a64:85b7ef5d4f9b9f10
% canonical-lessons: SA-C516-manoranjanam, SA-C516-hasyam, SA-C516-vinodayati, SA-C516-sangrahah, SA-C516-rucih

\chapter{Free Time}
\label{ch:sa-free-time}
\hlchaptermodality{516}

\begin{chapteropening}
\textbf{By the end of this chapter:} \emph{I can talk about what I do in my free time.}

Complete the last lesson of chapter 516: I can talk about what I do in my free time.
\end{chapteropening}

\section[\texorpdfstring{\sk{मनोरञ्जनम्} (manorañjanam)}{मनोरञ्जनम् (manorañjanam)}]{\sk{मनोरञ्जनम्} (manorañjanam) — entertainment}

\label{lesson:SA-C516-manoranjanam}



\begin{quote}
Before the new one: say the Sanskrit for ripe, then the Sanskrit for sucks.
\end{quote}

\subsection*{You'll want to know: \sk{मनोरञ्जनम्}}
\textbf{\sk{मनोरञ्जनम्}} — \emph{manorañjanam} — \textquotedblleft{}entertainment\textquotedblright{}.

\begin{quote}
\sk{रविवासरे} \sk{वयं} \sk{मनोरञ्जनाय} \sk{उद्यानं} \sk{गच्छामः।}
\end{quote}

\begin{quote}
\emph{ravivāsare vayaṃ manorañjanāya udyānaṃ gacchāmaḥ.}
\end{quote}

\begin{quote}
On Sunday we go to the park for some fun.
\end{quote}

A neuter noun.

\begin{cousinweb}
\textbf{\sk{मनोरञ्जनम्}} (\emph{manorañjanam}) is built from \emph{manas}, \textquotedblleft{}the mind\textquotedblright{}, and \emph{rañjana}, \textquotedblleft{}colouring, delighting\textquotedblright{}, from the root √rañj, \textquotedblleft{}to be coloured, to be pleased\textquotedblright{}; MW records \emph{manorañjana} only as the title of a book, so \textquotedblleft{}entertainment\textquotedblright{} is the modern sense (MW).
\end{cousinweb}

\begin{grammarlens}[title={What you've built}]
One more word for describing what you lived and felt.
\end{grammarlens}

\subsection*{Your turn}
\begin{itemize}
\raggedright
  \item \emph{Say it:} \emph{manorañjanam}
  \item \emph{Say it:} \emph{manorañjanam}, once more
  \item \emph{Say it:} \emph{ravivāsare vayaṃ manorañjanāya udyānaṃ gacchāmaḥ.} — On Sunday we go to the park for some fun.
  \item \emph{Recall:} say the Sanskrit for ripe, then the Sanskrit for sucks, then say \emph{manorañjanam} again
\end{itemize}

\begin{tcolorbox}[breakable,colback=teal!4,colframe=teal!35!black,title={Before you move on}]
What does \sk{मनोरञ्जनम्} mean? (\textquotedblleft{}Entertainment\textquotedblright{}.) Say it once more.
\end{tcolorbox}
\section[\texorpdfstring{\sk{हास्यम्} (hāsyam)}{हास्यम् (hāsyam)}]{\sk{हास्यम्} (hāsyam) — fun; laughter}

\label{lesson:SA-C516-hasyam}



\begin{quote}
Before the new one: say the Sanskrit for sucks, then the Sanskrit for entertainment.
\end{quote}

\subsection*{You'll want to know: \sk{हास्यम्}}
\textbf{\sk{हास्यम्}} — \emph{hāsyam} — \textquotedblleft{}fun; laughter\textquotedblright{}.

\begin{quote}
\sk{बालस्य} \sk{हास्यं} \sk{मह्यं} \sk{रोचते।}
\end{quote}

\begin{quote}
\emph{bālasya hāsyaṃ mahyaṃ rocate.}
\end{quote}

\begin{quote}
I love the child's laughter.
\end{quote}

A neuter noun.

\begin{cousinweb}
\textbf{\sk{हास्यम्}} (\emph{hāsyam}) comes from the root √has, \textquotedblleft{}to laugh\textquotedblright{}, which also gives \emph{hasati}; MW gives \emph{hāsya} as laughter, mirth, and jest, fun, amusement (MW).
\end{cousinweb}

\begin{grammarlens}[title={What you've built}]
One more word for describing what you lived and felt.
\end{grammarlens}

\subsection*{Your turn}
\begin{itemize}
\raggedright
  \item \emph{Say it:} \emph{hāsyam}
  \item \emph{Say it:} \emph{hāsyam}, once more
  \item \emph{Say it:} \emph{bālasya hāsyaṃ mahyaṃ rocate.} — I love the child's laughter.
  \item \emph{Recall:} say the Sanskrit for sucks, then the Sanskrit for entertainment, then say \emph{hāsyam} again
\end{itemize}

\begin{tcolorbox}[breakable,colback=teal!4,colframe=teal!35!black,title={Before you move on}]
What does \sk{हास्यम्} mean? (\textquotedblleft{}Fun; laughter\textquotedblright{}.) Say it once more.
\end{tcolorbox}
\section[\texorpdfstring{\sk{विनोदयति} (vinodayati)}{विनोदयति (vinodayati)}]{\sk{विनोदयति} (vinodayati) — amuses}

\label{lesson:SA-C516-vinodayati}



\begin{quote}
Before the new one: say the Sanskrit for entertainment, then the Sanskrit for fun.
\end{quote}

\subsection*{You'll want to know: \sk{विनोदयति}}
\textbf{\sk{विनोदयति}} — \emph{vinodayati} — \textquotedblleft{}amuses\textquotedblright{}.

\begin{quote}
\sk{पितामहः} \sk{गीतेन} \sk{बालान्} \sk{विनोदयति।}
\end{quote}

\begin{quote}
\emph{pitāmahaḥ gītena bālān vinodayati.}
\end{quote}

\begin{quote}
Grandfather amuses the children with a song.
\end{quote}

Past (he …): \sk{सः} \sk{बालान्} \sk{विनोदितवान्}, \emph{saḥ bālān vinoditavān}.

\begin{cousinweb}
\textbf{\sk{विनोदयति}} (\emph{vinodayati}) is built on \emph{vi}, \textquotedblleft{}away\textquotedblright{}, and the causative of the root √nud, \textquotedblleft{}to push\textquotedblright{}, which also gives \emph{nudati}: to drive dullness away; MW gives \emph{vinud} as to drive away, and to divert (MW).
\end{cousinweb}

\begin{grammarlens}[title={What you've built}]
One more word for describing what you lived and felt.
\end{grammarlens}

\subsection*{Your turn}
\begin{itemize}
\raggedright
  \item \emph{Say it:} \emph{vinodayati}
  \item \emph{Say it:} \emph{vinodayati}, once more
  \item \emph{Say it:} \emph{pitāmahaḥ gītena bālān vinodayati.} — Grandfather amuses the children with a song.
  \item \emph{Recall:} say the Sanskrit for entertainment, then the Sanskrit for fun, then say \emph{vinodayati} again
\end{itemize}

\begin{tcolorbox}[breakable,colback=teal!4,colframe=teal!35!black,title={Before you move on}]
What does \sk{विनोदयति} mean? (\textquotedblleft{}Amuses\textquotedblright{}.) Say it once more.
\end{tcolorbox}
\section[\texorpdfstring{\sk{संग्रहः} (saṅgrahaḥ)}{संग्रहः (saṅgrahaḥ)}]{\sk{संग्रहः} (saṅgrahaḥ) — a collection}

\label{lesson:SA-C516-sangrahah}



\begin{quote}
Before the new one: say the Sanskrit for fun, then the Sanskrit for amuses.
\end{quote}

\subsection*{You'll want to know: \sk{संग्रहः}}
\textbf{\sk{संग्रहः}} — \emph{saṅgrahaḥ} — \textquotedblleft{}a collection\textquotedblright{}.

\begin{quote}
\sk{एषः} \sk{मम} \sk{भ्रातुः} \sk{संग्रहः} \sk{अस्ति।}
\end{quote}

\begin{quote}
\emph{eṣaḥ mama bhrātuḥ saṅgrahaḥ asti.}
\end{quote}

\begin{quote}
This is my brother's collection.
\end{quote}

A masculine noun.

\begin{cousinweb}
\textbf{\sk{संग्रहः}} (\emph{saṅgrahaḥ}) is built on \emph{sam}, \textquotedblleft{}together\textquotedblright{}, and the root √grah, \textquotedblleft{}to seize, to take\textquotedblright{}; MW gives \emph{saṃgraha} as collecting, gathering, an accumulation (MW).
\end{cousinweb}

\begin{grammarlens}[title={What you've built}]
One more word for describing what you lived and felt.
\end{grammarlens}

\subsection*{Your turn}
\begin{itemize}
\raggedright
  \item \emph{Say it:} \emph{saṅgrahaḥ}
  \item \emph{Say it:} \emph{saṅgrahaḥ}, once more
  \item \emph{Say it:} \emph{eṣaḥ mama bhrātuḥ saṅgrahaḥ asti.} — This is my brother's collection.
  \item \emph{Recall:} say the Sanskrit for fun, then the Sanskrit for amuses, then say \emph{saṅgrahaḥ} again
\end{itemize}

\begin{tcolorbox}[breakable,colback=teal!4,colframe=teal!35!black,title={Before you move on}]
What does \sk{संग्रहः} mean? (\textquotedblleft{}A collection\textquotedblright{}.) Say it once more.
\end{tcolorbox}
\section[\texorpdfstring{\sk{रुचिः} (ruciḥ)}{रुचिः (ruciḥ)}]{\sk{रुचिः} (ruciḥ) — a liking}

\label{lesson:SA-C516-rucih}



\begin{quote}
Before the new one: say the Sanskrit for amuses, then the Sanskrit for a collection.
\end{quote}

\subsection*{You'll want to know: \sk{रुचिः}}
\textbf{\sk{रुचिः}} — \emph{ruciḥ} — \textquotedblleft{}a liking\textquotedblright{}.

\begin{quote}
\sk{मम} \sk{भगिन्याः} \sk{संगीते} \sk{रुचिः} \sk{अस्ति।}
\end{quote}

\begin{quote}
\emph{mama bhaginyāḥ saṅgīte ruciḥ asti.}
\end{quote}

\begin{quote}
My sister has a liking for music.
\end{quote}

A feminine noun.

\begin{cousinweb}
\textbf{\sk{रुचिः}} (\emph{ruciḥ}) comes from the root √ruc, \textquotedblleft{}to shine, to be pleasing\textquotedblright{}, which also gives \emph{rocate}; MW gives \emph{ruci} as light, and as liking, taste, relish (MW).
\end{cousinweb}

\begin{grammarlens}[title={What you've built}]
One more word for describing what you lived and felt.
\end{grammarlens}

\subsection*{Your turn}
\begin{itemize}
\raggedright
  \item \emph{Say it:} \emph{ruciḥ}
  \item \emph{Say it:} \emph{ruciḥ}, once more
  \item \emph{Say it:} \emph{mama bhaginyāḥ saṅgīte ruciḥ asti.} — My sister has a liking for music.
  \item \emph{Recall:} say the Sanskrit for amuses, then the Sanskrit for a collection, then say \emph{ruciḥ} again
\end{itemize}

\begin{tcolorbox}[breakable,colback=teal!4,colframe=teal!35!black,title={Before you move on}]
What does \sk{रुचिः} mean? (\textquotedblleft{}A liking\textquotedblright{}.) Say it once more.
\end{tcolorbox}
